% Страница «Как пользоваться шпорой» — идёт сразу после оглавления.
\clearpage
\renewcommand\pfKolontitul{Как пользоваться шпорой}
\pdfbookmark[1]{Как пользоваться шпорой}{bm-kak}
\pfZagolovok{Как пользоваться шпорой}

\opr{Одна тема — одна плашка.} У каждой темы есть код, например \pfChip{Ч1.3}.
Буква — раздел, первая цифра — глава, вторая — номер темы в главе.

\begin{pfTablica}[bez shapki]{>{\color{pfBrend}\pfBrendZhirnyj}cl>{\color{pfBrend}\pfBrendZhirnyj}cl}
Ч & числа и вычисления  & П & планиметрия\\
А & алгебра             & С & стереометрия\\
Т & тригонометрия       & К & векторы и координаты\\
Ф & функции и графики   & В & вероятность и статистика\\
Д & производная         & М & текстовые и экономические задачи\\
\end{pfTablica}

\smallskip
\opr{Внизу плашки — связи.} «Опирается на» — что стоит повторить сначала.
«Где встречается» — типы заданий ЕГЭ, в которых тема нужна, со ссылкой на разборы.
«Задания ЕГЭ» — номера заданий, где тема пригодится.

\begin{primer}
синяя полоска — короткий пример применения.
\end{primer}
\begin{oshibka}
оранжевая — то, на чём чаще всего теряют баллы.
\end{oshibka}

\opr{Всё кликается.} Код темы ведёт к теме, «Оглавление$\,\uparrow$» внизу страницы —
к оглавлению, «Задания ЕГЭ$\,\uparrow$» — к списку заданий: по каждому номеру там собраны
все нужные темы со страницами. Фиолетовые номера в строке «Задания ЕГЭ» под темой ведут к разборам этого задания в практике. Синяя вкладка на правом краю показывает раздел и ведёт к началу главы.

\opr{Практика} — разборы реальных прототипов по номерам заданий ЕГЭ. Каждое задание
разбито на типы, как в банке задач Профиматики: у каждого типа есть ссылка на задачи
с разборами в банке. В начале разбора написано, какие темы понадобятся. А в самом конце книги —
раздел «Задания ЕГЭ $\to$ темы»: какие темы повторить перед конкретным номером.
